\documentclass[11pt]{article}
\usepackage{mathblatt}
% gesetzt von werkzeuge/setzer.py (Sorte tisch) aus dem Lernweg FUR
% und den Bankzeilen; Muster bau/proben/2026-10-09/hypotenuse-muster.
\usepackage{enumitem}
\usepackage{adjustbox,varwidth}
\newgeometry{a4paper,margin=18mm,top=15mm,bottom=20mm,footskip=9mm}
\pagestyle{fancy}\fancyhf{}\renewcommand{\headrulewidth}{0pt}
\fancyfoot[L]{\footnotesize\color{mbgrau}FUR \textperiodcentered{} Lösungen \textperiodcentered{} Zahlen vergleichen – auch negativ, mit Bruch, Prozent, Quadrat und Wurzel \textperiodcentered{} Tischblatt}
\fancyfoot[R]{\footnotesize\color{mbgrau}\thepage}
\setlength{\parskip}{3pt}
\renewcommand{\kreuz}[1]{\mbox{$\square$\ #1}\quad}
\providecommand{\kk}[1]{$\square$\,#1\hspace{0.9em}}
\providecommand{\linie}[1]{\,{\color{black!45}\rule[-1pt]{#1}{0.4pt}}\,}
\providecommand{\grau}[1]{{\color{mbgrau}#1}}
\newcommand{\kreuzl}[1]{\par\hangindent1.4em\hangafter1\noindent$\square$\ #1}
\newcommand{\aufg}[3]{\par\noindent\makebox[0pt][r]{#3}\begin{minipage}[t]{\linewidth}#1\end{minipage}\par}
\newcommand{\aufgg}[4]{\par\noindent\makebox[0pt][r]{#4}\begin{minipage}[t]{0.6\linewidth}\vspace{0pt}#1\end{minipage}\hfill
  \begin{minipage}[t]{0.37\linewidth}\vspace{0pt}\centering #2\end{minipage}\par}
\newcommand{\nr}[1]{\makebox[7mm][l]{\textbf{#1}}}
\newcommand{\lsg}[3]{\par\noindent\begin{minipage}[t]{9mm}\textbf{#1}\end{minipage}%
  \begin{minipage}[t]{52mm}\raggedright\bfseries\boldmath #2\end{minipage}\hspace{3mm}%
  \begin{minipage}[t]{\dimexpr\linewidth-64mm\relax}\raggedright\small #3\end{minipage}\par\vspace{4pt}}
\newlength{\kr}
\makeatletter
\newcommand{\karomerk}[2]{\protected@write\@auxout{}{\string\karoseite{#1}{#2}{\thepage}}}
\newcommand{\karoseite}[3]{\expandafter\xdef\csname karo@#1@#2\endcsname{#3}}
\newcommand{\karohinweis}[3]{\karomerk{h}{#1}\ifcsname karo@k@#1\endcsname
  \ifnum\csname karo@k@#1\endcsname=0\csname karo@h@#1\endcsname\relax#2\else#3\fi
  \else#3\fi}
\makeatother
\begin{document}

\noindent{\Large\bfseries Lösungen}\hspace{1em}{\color{mbgrau}Zahlen vergleichen – auch negativ, mit Bruch, Prozent, Quadrat und Wurzel}\par\smallskip
\noindent{\small Links das Ergebnis, rechts der Weg in kurzen Schritten. Stimmt dein Ergebnis nicht, such den ersten Schritt, der anders ist.}\par
\par\Needspace{25mm}\vspace{6pt}\noindent\colorbox{black!12}{\makebox[7mm]{\bfseries\strut A}}\hspace{2mm}{\bfseries Dezimalzahlen vergleichen}\par\vspace{4pt}
\lsg{A1a)}{Zehntel}{}
\lsg{b)}{Tausendstel}{}
\lsg{c)}{Einer}{}
\lsg{d)}{Hundertstel}{}
\lsg{A2a)}{$2{,}4 > 2{,}38$}{\ (mit Nullen: $2{,}40$ und $2{,}38$)}
\lsg{b)}{$0{,}09 < 0{,}1$}{\ (mit Nullen: $0{,}09$ und $0{,}10$)}
\lsg{c)}{$5{,}05 < 5{,}5$}{\ (mit Nullen: $5{,}05$ und $5{,}50$)}
\lsg{d)}{$1{,}250 = 1{,}25$}{\ (mit Nullen: $1{,}250$ und $1{,}250$)}
\lsg{A3}{$0{,}06 < 0{,}066 < 0{,}6 < 0{,}61$}{$0{,}06$; $0{,}066$; $0{,}6$; $0{,}61$ – mit Nullen: $0{,}060 < 0{,}066 < 0{,}600 < 0{,}610$}
\par\Needspace{25mm}\vspace{6pt}\noindent\colorbox{black!12}{\makebox[7mm]{\bfseries\strut B}}\hspace{2mm}{\bfseries Negative Zahlen vergleichen}\par\vspace{4pt}
\lsg{B1a)}{$-2$}{}
\lsg{b)}{$-0{,}5$}{}
\lsg{c)}{$-3$}{}
\lsg{d)}{$-0{,}1$}{}
\lsg{B2a)}{$-1{,}2 < -1{,}02$}{}
\lsg{b)}{$-0{,}3 < 0{,}1$}{}
\lsg{c)}{$-2{,}5 < -2{,}45$}{}
\lsg{d)}{$-0{,}09 > -0{,}1$}{}
\lsg{B3a)}{$-0{,}3$}{Abstand $0{,}2$, Hälfte $0{,}1$; $-0{,}4 + 0{,}1 = -0{,}3$}
\lsg{b)}{$-0{,}8$}{Abstand $0{,}4$, Hälfte $0{,}2$; $-1 + 0{,}2 = -0{,}8$}
\lsg{c)}{$-0{,}1$}{Abstand $0{,}4$, Hälfte $0{,}2$; $-0{,}3 + 0{,}2 = -0{,}1$}
\lsg{B4}{Kiew, Oslo, Riga, Prag, Wien}{$-4$; $-3{,}5$; $-3{,}05$; $-0{,}8$; $0{,}5$ – also Kiew, Oslo, Riga, Prag, Wien; am kältesten in Kiew.}
\par\Needspace{25mm}\vspace{6pt}\noindent\colorbox{black!12}{\makebox[7mm]{\bfseries\strut C}}\hspace{2mm}{\bfseries Bruch und Prozent: erst umwandeln}\par\vspace{4pt}
\lsg{C1a)}{$0{,}75$}{$\frac{3}{4} = \frac{75}{100} = 0{,}75$}
\lsg{b)}{$0{,}35$}{$\frac{7}{20} = \frac{35}{100} = 0{,}35$}
\lsg{c)}{$0{,}125$}{$\frac{1}{8} = \frac{125}{1000} = 0{,}125$}
\lsg{d)}{$0{,}36$}{$\frac{9}{25} = \frac{36}{100} = 0{,}36$}
\lsg{C2a)}{$0{,}45$}{$45\,\% = 0{,}45$}
\lsg{b)}{$0{,}045$}{$4{,}5\,\% = 0{,}045$}
\lsg{c)}{$1{,}2$}{$120\,\% = 1{,}2$}
\lsg{d)}{$0{,}005$}{$0{,}5\,\% = 0{,}005$}
\lsg{C3a)}{$6\,\% < 0{,}6$}{$0{,}06$ und $0{,}6$, also $6\,\% < 0{,}6$}
\lsg{b)}{$0{,}35 = \frac{7}{20}$}{$0{,}35$ und $0{,}35$, also $0{,}35 = \frac{7}{20}$}
\lsg{c)}{$\frac{2}{5} < 45\,\%$}{$0{,}4$ und $0{,}45$, also $\frac{2}{5} < 45\,\%$}
\lsg{d)}{$\frac{1}{8} < 0{,}13$}{$0{,}125$ und $0{,}13$, also $\frac{1}{8} < 0{,}13$}
\par\Needspace{25mm}\vspace{6pt}\noindent\colorbox{black!12}{\makebox[7mm]{\bfseries\strut D}}\hspace{2mm}{\bfseries Quadrat und Wurzel}\par\vspace{4pt}
\lsg{D1a)}{$0{,}04$}{$0{,}2^2 = 0{,}2 \cdot 0{,}2 = 0{,}04$}
\lsg{b)}{$0{,}25$}{$0{,}5^2 = 0{,}5 \cdot 0{,}5 = 0{,}25$}
\lsg{c)}{$1{,}21$}{$1{,}1^2 = 1{,}1 \cdot 1{,}1 = 1{,}21$}
\lsg{d)}{$0{,}01$}{$0{,}1^2 = 0{,}1 \cdot 0{,}1 = 0{,}01$}
\lsg{D2a)}{$0{,}3$}{$\sqrt{0{,}09} = 0{,}3$, denn $0{,}3 \cdot 0{,}3 = 0{,}09$}
\lsg{b)}{$1{,}2$}{$\sqrt{1{,}44} = 1{,}2$, denn $1{,}2 \cdot 1{,}2 = 1{,}44$}
\lsg{c)}{$0{,}2$}{$\sqrt{0{,}04} = 0{,}2$, denn $0{,}2 \cdot 0{,}2 = 0{,}04$}
\lsg{d)}{$0{,}7$}{$\sqrt{0{,}49} = 0{,}7$, denn $0{,}7 \cdot 0{,}7 = 0{,}49$}
\lsg{D3a)}{$1{,}7 < \sqrt{3} < 1{,}8$}{$1{,}7^2 = 2{,}89 < 3 < 3{,}24 = 1{,}8^2$, also $1{,}7 < \sqrt{3} < 1{,}8$}
\lsg{b)}{$2{,}2 < \sqrt{5} < 2{,}3$}{$2{,}2^2 = 4{,}84 < 5 < 5{,}29 = 2{,}3^2$, also $2{,}2 < \sqrt{5} < 2{,}3$}
\lsg{c)}{$0{,}7 < \sqrt{0{,}5} < 0{,}8$}{$0{,}7^2 = 0{,}49 < 0{,}5 < 0{,}64 = 0{,}8^2$, also $0{,}7 < \sqrt{0{,}5} < 0{,}8$}
\par\Needspace{25mm}\vspace{6pt}\noindent\colorbox{black!12}{\makebox[7mm]{\bfseries\strut E}}\hspace{2mm}{\bfseries Alles gemischt: vergleichen und entscheiden}\par\vspace{4pt}
\lsg{E1}{$8\,\% = 0{,}08$}{$8\,\% = 0{,}08$; die anderen sind $0{,}9$, $0{,}09$ und $0{,}085$.}
\lsg{E2}{$-\frac{3}{4} < -0{,}7 < 0{,}3 < \sqrt{0{,}16}$}{$-0{,}75$; $-0{,}7$; $0{,}3$; $0{,}4$ – also $-\frac{3}{4} < -0{,}7 < 0{,}3 < \sqrt{0{,}16}$}
\lsg{E3}{$\frac{5}{8} > 0{,}6$}{$\frac{5}{8} = 0{,}625 > 0{,}600$. Falsch: $\sqrt{2} \approx 1{,}41 < 1{,}5$; $0{,}3^2 = 0{,}09 < 0{,}3$; $-0{,}5 > -0{,}6$.}
\lsg{E4}{Laden A; 3 Prozentpunkte mehr als B}{Laden A: ein Viertel $= \frac{1}{4} = 25\,\%$; B: $22\,\%$; C: jeder fünfte Euro $= \frac{1}{5} = 20\,\%$. Am größten bei A, um $25 - 22 = 3$ Prozentpunkte mehr als bei B.}
\lsg{E5}{Ja; Fisch in $3{,}6$\,m Tiefe}{Tiefen als Zahlen: $-4{,}8$ und $-2{,}4$; Abstand $2{,}4$, Hälfte $1{,}2$; $-4{,}8 + 1{,}2 = -3{,}6$. Ja, der Fisch ist in $3{,}6$\,m Tiefe, das ist tiefer als $3{,}5$\,m.}
\end{document}
